\documentclass[11pt]{article}

\usepackage[margin=1in]{geometry}
\usepackage{amsmath, amssymb}
\usepackage{booktabs}
\usepackage[hidelinks]{hyperref}
\usepackage[htt]{hyphenat}
\sloppy

\title{PINO and PI-DeepONet baselines: what the implementation actually does}
\date{\today}

\begin{document}
\maketitle

\noindent
Summary of the two published label-free operator-learning baselines on branch
\texttt{baselines-pino-deeponet}, read off the source rather than the papers. They exist so the
comparison is against \emph{published} physics-informed operator training and not only against our
own $t=0$ negative control. Code lives in \texttt{tensorpils/baselines/}; nothing in
\texttt{losses.py}, \texttt{trainer.py} or \texttt{physics.py} is modified, and only
\texttt{cli.py} gains dispatch.

\section*{The short answers}

\begin{center}
\begin{tabular}{lll}
\toprule
 & \textbf{PINO} & \textbf{PI-DeepONet} \\
\midrule
Architecture      & FNO (ours, unchanged)          & DeepONet (new, unstacked) \\
Derivatives       & 2nd-order central differences  & autodiff through the trunk \\
Dirichlet BC      & hard, boundary nodes zeroed & hard, $\sin(\pi x)\sin(\pi y)$; or zeroed + soft penalty (\texttt{--pideeponet\_bc zero}) \\
Residual form     & strong                         & strong \\
Evaluated at      & grid nodes, interior only      & collocation points \\
Reduction         & relative $L^2$ ratio           & relative $L^2$ ratio \\
\bottomrule
\end{tabular}
\end{center}

\paragraph{Does PINO just sit on top of the FNO?} Yes. \texttt{ZeroBoundaryModel} wraps our
\texttt{FNOModel} unchanged and zeroes its boundary nodes; the loss is swapped and nothing else.
The only architectural addition in the whole branch is the DeepONet.

\section*{Boundary conditions}

Both baselines impose $u=0$ on $\partial\Omega$ \textbf{hard}, as a property of the \emph{model}
rather than of the loss, so training and evaluation see the same function. \textbf{They do so by
different mechanisms, and deliberately so.}

\paragraph{PINO --- zeroing the boundary nodes (default).} The prediction is multiplied by a
0/1 mask that is zero on the outer ring of nodes. This is exactly the operation
\texttt{losses.py} applies to our own arms via \texttt{apply\_zero\_boundary}, so PINO and the
\texttt{galerkin}/\texttt{pls} arms now impose the Dirichlet condition identically and differ only
in the residual --- strong versus weak form, finite differences versus exact $Q_1$ assembly.
A unit test asserts the two operations agree bit-for-bit.

This departs from the reference implementation, which multiplies by
$\sin(\pi x)\sin(\pi y)$ (still available under \texttt{--pino\_bc mollifier}). The reason is that
the mollifier injects a smooth global decay profile into the model, and for the multi-frequency
sine dataset used here that profile is close to the solution's own structure --- information our
arms never receive. It also does not generalise: zeroing boundary nodes works on any mesh, whereas
a mollifier must be hand-built per geometry. \emph{Expect this to cost PINO accuracy relative to
the mollified version}; that is the price of making the comparison attributable.

Masking is bit-exact, so unlike the mollifier it needs no boundary snap: $\sin(\pi\cdot 1.0)$
evaluates to $8.7\times10^{-8}$ in \texttt{float32}, not $0$.

\paragraph{PI-DeepONet --- the mollifier (default), or zeroing plus a penalty.} By default
unchanged from the reference: $\sin(\pi x)\sin(\pi y)$, evaluated at the \emph{query}
coordinates inside \texttt{forward\_at}, so the autodiff Laplacian differentiates the mollified
field as it must. The same objection applies here, so \texttt{--pideeponet\_bc zero} provides the
counterpart of \texttt{--pino\_bc zero}: the grid output's boundary nodes are zeroed exactly as our
arms do, and the residual gets the original PI-DeepONet soft boundary penalty, weighted by
\texttt{--pi\_lambda\_bc}. The penalty is not optional there. PINO's finite-difference stencil
touches the zeroed boundary ring, so masking alone couples the interior residual to the BC; a
pointwise autodiff Laplacian at interior points is blind to the boundary values and to any mask on
them, and without the penalty the objective is invariant under adding harmonic functions. To make
$\lambda_{\mathrm{bc}}=1$ meaningful at any solution magnitude the penalty is written in the
residual's own units: per sample, the boundary RMS divided by the (detached) interior RMS of the
prediction, averaged over the batch --- dimensionless like the relative residual
$\|Lu-f\|/\|f\|$. This matters: at $K=10$ the solutions have RMS $\approx1.6\cdot10^{-3}$ while
the relative residual is $O(1)$, so a plain mean-square boundary term at unit weight would be five
to six orders of magnitude too weak. The \texttt{mse} reduction pairs with the plain mean-square boundary penalty of
the original paper instead.

\textbf{This is an advantage the baselines have over our bare \texttt{galerkin} arm}, which does not
get a hard BC. It is deliberate: skipping it would be strawmanning. The soft BC penalty of the
original PI-DeepONet ($\lambda_{\mathrm{bc}}$) is therefore $0$ by default, which removes a tuned
hyperparameter rather than adding one.

\section*{Stokes: the saddle-point port}

Both baselines exist for Stokes as well (\texttt{--pde stokes --loss pino|pideeponet}). There is
no reference implementation for a saddle point, so the port is the Darcy recipe applied to each
equation of $-\mu\Delta u + \nabla p = f$, $\nabla\cdot u = 0$, with three decisions that are the
Stokes form of choices already fixed above. (i) The residual is evaluated on the velocity grid at
interior nodes, and the \emph{pressure is taken on the full fine grid} (the FNO's third channel
times \texttt{p\_scale}) rather than on the $[::2,::2]$ Q1 subgrid the FEM arms read: a strong form
needs $\nabla p$ at every interior velocity node. Evaluation is unchanged. (ii) The two equations
are \emph{stacked} in one relative residual, $\|(r_{\mathrm{mom}}, w\,r_{\mathrm{div}})\|/\|(f,0)\|$
(\texttt{stokes\_reduce}), the strong-form analogue of the negative control $\tfrac12\|Kc-b\|^2$ with
$b=(M_u f, 0)$. The continuity weight $w$ (\texttt{--pi\_div\_weight}) is the one hyperparameter the
saddle point adds --- $\nabla\cdot u\sim k u$ against $f\sim\mu k^2 u$, so the rows have different
units --- and is tuned on validation like a learning rate. (iii) The hard BC, by either mechanism,
acts on the velocity pair only (\texttt{DeepONetModel(bc\_channels=(0,1))}); the pressure has a gauge,
not a boundary condition, fixed by projection at evaluation and invisible to $\nabla p$. For PINO no
wrapper is needed at all: the trainer zeroes the velocity boundary ring itself, exactly as
\texttt{\_predict} projects every arm. The DeepONet emits three channels through $C$ independent
groups of $p$ basis functions sharing the hidden layers. Tests pin the port: the FD residual of a
manufactured Stokes solution converges at rate $2$ (which also fixes the axis convention), the
autodiff operators match the analytic derivatives, and the FD and autodiff losses agree on a resolved
field. No Stokes baseline results exist yet; \texttt{experiments/baselines/stokes/} holds the sweep.

\section*{How derivatives are computed}

\paragraph{PINO: finite differences.} The Laplacian comes from
\texttt{neuralop.losses.differentiation.FiniteDiff} --- the reference library's own utility,
second-order central in the interior, pure \texttt{torch} so it stays differentiable. Using the
library's function rather than a reimplementation removes ``you coded the baseline wrong'' as a
reading of any result.

The choice of scheme is not taste. The reference repo uses FFT differentiation for its
\emph{periodic} problems (Burgers, Navier--Stokes) and central differences for \emph{Dirichlet}
Darcy. Every PDE here is homogeneous Dirichlet on the unit square, so Darcy is the analogue and
finite differences is the faithful choice.

Note \texttt{FiniteDiff} names axis $-2$ ``x'' while this repo names it $y$; the spacings are passed
in its order, $(h_y, h_x)$.

\paragraph{PI-DeepONet: autodiff.} The trunk consumes the coordinate $y=(x,y)$ directly, so the
Laplacian is exact: one \texttt{grad} of $u$ with respect to \texttt{coords}, then a second
\texttt{grad} per dimension, with \texttt{create\_graph=True} throughout so the residual remains
differentiable in the parameters.

\textbf{Coordinates are passed per sample}, as \texttt{[B, Q, 2]} rather than \texttt{[Q, 2]}. With
the shared form, \texttt{grad(u.sum(), coords)} returns $\sum_b \partial u_b/\partial y$ --- the
batch \emph{sum} of gradients --- and the per-sample residual cannot be recovered. The code raises
if given the shared form.

\section*{The DeepONet architecture}

Standard unstacked DeepONet,
\begin{equation*}
   \mathcal G_\theta(f)(y) = \frac{1}{\sqrt{p}}\sum_{k=1}^{p} b_k(f)\,t_k(y) + b_0 ,
\end{equation*}
with an MLP branch over the input function sampled on the sensor grid and an MLP trunk over the
query coordinate. \texttt{neuralop} ships no DeepONet, so this is a direct implementation.

\begin{center}
\begin{tabular}{ll}
\toprule
\textbf{Component} & \textbf{Value} \\
\midrule
Basis functions $p$      & $128$ \\
Branch                   & MLP $4225 \to 256 \to 256 \to 256 \to 128$, $\tanh$ \\
Trunk                    & MLP $128 \to 256 \to 256 \to 256 \to 128$, $\tanh$ \\
Trunk input              & $64$ random Fourier features (fixed, scale $2.0$) \\
Sensors                  & the full $65^2$ grid, flattened \\
Trainable parameters     & $1.44\times10^{6}$ (branch $1.25\times10^{6}$, trunk $0.20\times10^{6}$) \\
\bottomrule
\end{tabular}
\end{center}

Three implementation choices exist so that a poor result is attributable to the objective and not to
a handicapped baseline:
\begin{itemize}
   \item \textbf{Random Fourier features on the trunk.} A plain MLP has a strong spectral bias and
   these data are multi-frequency ($K=4$ means modes up to $4\times4$); without them the DeepONet
   underfits for reasons unrelated to the loss. Standard practice in the PI-DeepONet literature.
   The matrix $B$ is a fixed buffer, so it travels with the checkpoint.
   \item \textbf{$1/\sqrt{p}$ output scaling.} A dot product of two $p$-vectors grows like
   $\sqrt{p}$, so without this the output magnitude at initialisation depends on $p$ (measured
   $5.5\times$ the target at $p=32$, $12.4\times$ at $p=128$). Starting an operator an order of
   magnitude too large costs the baseline real epochs.
   \item \textbf{Branch input scaling} \texttt{f\_scale}: an MLP on a raw $O(10^2)$ input trains
   badly, while the FNO's lifting layer absorbs the scale by itself.
\end{itemize}

\paragraph{It is also a drop-in for the FNO.} \texttt{forward(f\_grid)} returns
\texttt{[B,1,H,W]}, the \texttt{FNOModel} signature, so the DeepONet can be paired with
\emph{any} loss in \texttt{losses.py} through the unmodified \texttt{PoissonTrainer}. This fills the
``our loss, other architecture'' cell for free. \texttt{forward\_at(f\_grid, coords)} is the second
entry point, and is what the autodiff residual uses.

\section*{Things that are easy to miss}

\begin{itemize}
   \item \textbf{The reduction is a relative ratio, not our sum-over-nodes convention.} PINO reduces
   with $\mathrm{mean}_b\,\|{-\Delta_h u} - f\|_2 / \|f\|_2$. It is inherited, not a physics-informed design choice: this is \texttt{LpLoss.rel()} from the reference FNO/PINO codebase, the relative-$L^2$ reduction introduced as the \emph{supervised} training loss in the original FNO paper and reused unchanged for the equation residual. It changes the gradient scale by
   \emph{orders of magnitude} relative to our $\tfrac12\|\cdot\|^2$ losses, which is why each
   baseline arm gets its own learning-rate sweep. Comparing at a shared learning rate would be
   meaningless. \texttt{--pino\_reduction mse} puts them back on our footing when that is the
   question.
   \item \textbf{It is a ratio of norms, not of squared norms}, and that is not cosmetic:
   $\partial_r \big(\|r\|_2/\|f\|_2\big) = r/(\|r\|_2\|f\|_2)$ has norm $1/\|f\|_2$ for
   \emph{every} $r\neq0$, so the gradient magnitude does not vanish as the residual does (measured
   constant across six orders of magnitude in $\|r\|$, while $\tfrac12\|r\|^2$ decays linearly).
   PINO's objective is therefore not a smooth least-squares problem near its solution, its floor is
   set by the learning rate, and the Hessian analysis of Section~2.3---which presumes the squared
   form---does not transfer to it as written.
   \item \textbf{The residual is evaluated on interior points only}, boundary rows and columns
   sliced off, as the reference does.
   \item \textbf{Allen--Cahn: strong vs.\ weak, same time discretisation.} Where our FEM residual is
   $M\frac{u^{n+1}-u^n}{\tau} + a^2Au^{n+1} - M\rho$, PINO drops the mass matrix and replaces
   $A \leftrightarrow -\Delta_h$. The time stepping (convex--concave or backward Euler) is identical,
   so the only thing under test is weak versus strong form.
   \item \textbf{Allen--Cahn reduces with \texttt{mse}, not the relative ratio}: the step residual
   has no right-hand side to normalise against.
   \item \textbf{Collocation defaults to the grid nodes} --- for fairness (every arm then sees the
   same points), and for Allen--Cahn out of necessity: the residual needs $u^n$ at each collocation
   point and the FEM Newton reference exists only at nodes. Off-node collocation would require
   interpolating the reference and would inject an error unrelated to the method.
   For Poisson the source is closed-form, so random collocation is also available.
   \item \textbf{\texttt{detach\_coupling} mirrors our \texttt{ACGalerkinLoss}} and is set from
   \texttt{--bptt\_mode} exactly as the FEM arm is, so the two differ in the residual and not in
   what the gradient flows through.
   \item \textbf{Scoring is unchanged}: both are evaluated with the repo's FEM relative-$L^2$ metric
   and the eval-time boundary projection, so their numbers drop straight into the existing tables.
\end{itemize}

\section*{Invocation}

\begin{verbatim}
--loss pino                        # FNO + FD strong form (boundary nodes zeroed)
--loss pideeponet --model deeponet # DeepONet + autodiff strong form
--model deeponet --loss pls        # our loss, other architecture
\end{verbatim}
Relevant flags: \texttt{--mollify \{auto,on,off\}} (auto = on for the baseline losses, off
otherwise, so existing arms are unchanged), \texttt{--pino\_bc \{zero,mollifier\}} (default
\texttt{zero}), \texttt{--pideeponet\_bc \{mollifier,zero\}} (default \texttt{mollifier}) with
\texttt{--pi\_lambda\_bc}, \texttt{--pi\_n\_colloc}, \texttt{--pino\_reduction},
\texttt{--deeponet\_p/\_width/\_depth}, \texttt{--trunk\_fourier[\_scale]}.

\end{document}
